\chapter{Glossary}

\begin{description}
\item[Accept] Create a connected socket from a listening socket's pending connection.
\item[Address] A network-layer locator associated with an interface and routing context.
\item[Application server] A process that hosts application code behind a defined interface.
\item[ASGI] Asynchronous Server Gateway Interface for Python applications and servers.
\item[Backlog] An operating-system limit/hint for pending connections to a listener.
\item[Backpressure] Mechanisms that prevent producers from overwhelming finite consumers.
\item[Bind] Associate a socket with a local address and port.
\item[Body] HTTP/1.1 message bytes after the field-section delimiter, as framed by rules.
\item[Byte] Eight bits; the unit exposed by Python socket reads and writes.
\item[Cache] Component that stores responses for protocol-governed reuse.
\item[Client] The role that initiates a request in an interaction.
\item[Connection] Stateful transport relationship between endpoints.
\item[Content] HTTP message data after transfer coding is removed.
\item[Content negotiation] Selecting a representation using request preferences and availability.
\item[CRLF] Carriage return followed by line feed, bytes 13 and 10.
\item[CRUD] Create, read, update, delete; a data-operation mnemonic, not an HTTP specification.
\item[Datagram] A message-oriented transport unit, as exposed by UDP.
\item[DNS] Distributed hierarchical Domain Name System.
\item[Endpoint] An address and transport port, plus protocol context.
\item[Encoding] Mapping between abstract information such as text and bytes.
\item[Field] Named HTTP metadata with field-specific syntax and semantics.
\item[Flow control] Sender control intended to avoid overwhelming the receiver.
\item[Forward proxy] Intermediary selected on behalf of a client.
\item[Frame] A bounded protocol unit; often specifically a link-layer unit.
\item[Framing] Rules that locate message boundaries in bytes.
\item[Gateway] Intermediary that fronts or translates another system.
\item[Handler] Application callable that turns a Request into a Response.
\item[Header section] Start-line-following HTTP field lines ending at an empty line.
\item[Host] HTTP/1.1 field identifying target authority.
\item[HTTP] Stateless application-level request/response protocol family.
\item[HTTPS] HTTP whose communication is secured for the HTTPS origin rules, commonly with TLS.
\item[Idempotent] Having the same intended effect when repeated as when performed once.
\item[Intermediary] Proxy, gateway, cache, or tunnel participant between endpoints.
\item[IP] Internet Protocol, responsible for addressed packet forwarding.
\item[JSON] Text data-interchange syntax used for Task representations.
\item[Keep-alive] Reuse of a connection for more than one HTTP exchange.
\item[Latency] Elapsed delay across a defined operation or path.
\item[Listener] Socket configured to accept incoming transport connections.
\item[Loopback] Host-internal network path, represented by 127.0.0.1 or ::1.
\item[Media type] Identifier describing the format/nature of content.
\item[Method] Case-sensitive HTTP token expressing request semantics.
\item[Middleware] Callable that wraps downstream request handling.
\item[NAT] Translation of network addresses and often ports across routing contexts.
\item[Origin] Scheme, host, and port tuple governing an HTTP resource namespace.
\item[Packet] Network-layer unit forwarded by IP.
\item[Parameterization] Supplying SQL values separately from statement structure.
\item[Parser] Component that validates and transforms wire syntax into structured state.
\item[Port] Sixteen-bit transport identifier used in TCP/UDP endpoints.
\item[Problem details] Structured HTTP error representation standardized by RFC 9457.
\item[Protocol] Shared rules for communication between independent participants.
\item[Representation] Transferable current or intended state of a resource plus metadata.
\item[Resource] Conceptual target identified by a URI.
\item[REST] Architectural style with constraints including uniform interface and statelessness.
\item[Reverse proxy] Intermediary accepting traffic on behalf of origin applications.
\item[Route] Mapping from method/path shape to a handler.
\item[Router] Component that resolves requests against registered routes.
\item[Safe method] Method whose defined semantics are read-only.
\item[Segment] Common name for a TCP protocol data unit; also a path component by context.
\item[Server] Role that accepts requests and produces responses.
\item[Socket] Operating-system interface/object for network communication.
\item[Stateless] Requiring each request to carry the context needed for its interpretation.
\item[Status code] Three-digit HTTP response result code.
\item[Stream] Ordered sequence of bytes without application message boundaries.
\item[TCP] Reliable ordered byte-stream transport protocol.
\item[TLS] Protocol providing protected channels and authenticated identity under verification.
\item[Transaction] Group of database operations committed atomically or rolled back.
\item[URI] Uniform Resource Identifier.
\item[URL] URI that identifies a resource through a scheme with access semantics.
\item[UTF-8] Variable-length Unicode encoding compatible with ASCII for its first 128 values.
\item[Validation] Checking data against the grammar and rules owned by a layer.
\item[WAL] SQLite write-ahead logging mode.
\item[Worker] Execution unit handling queued connections or requests.
\item[WSGI] Synchronous Python Web Server Gateway Interface.
\end{description}
